\chapter{TEOREMA GELFAND-NAIMARK}\label{Gelfand_Naimark}

Seperti pada Bab~\ref{spectrum}, dalam bab ini---dan untuk sisa catatan ini---kita mengasumsikan bahwa medan
  \index{konvensi!dalam bab~\ref{Gelfand_Naimark} dan bab-bab berikutnya skalar bersifat kompleks}%
skalar kita adalah medan $\C$ bilangan kompleks.



\section{Ideal-Ideal Maksimal dalam $\fml C(X)$}
\begin{prop}\label{0006801} Jika $J$ adalah ideal sejati dalam suatu aljabar Banach beridentitas, maka
tutupannya juga merupakan ideal sejati.
\end{prop}

\begin{cor} Setiap ideal maksimal dalam aljabar Banach beridentitas bersifat tertutup.
\end{cor}

\begin{exam} Untuk setiap subhimpunan $C$ dari ruang topologis~$X$, himpunan
 \index{ideal!dalam $\fml C(X)$}%
 \index{J@$J_C$ (fungsi kontinu yang lenyap pada~$C$)}%
   \[ J_C := \bigl\{f \in \fml C(X)\colon f^{\sto}(C) = \{0\}\,\bigr\} \]
merupakan ideal dalam $\fml C(X)$. Lebih lanjut, $J_C \supseteq J_D$ apabila
$C~\subseteq~D~\subseteq~X$. (Selanjutnya kita akan menulis $J_x$ untuk ideal
$J_{\{x\}}$.)
\end{exam}

\begin{prop}\label{0006832} Misalkan $X$ ruang topologis kompak dan $I$ ideal sejati
dalam $\fml C(X)$. Maka terdapat $x \in X$ sedemikian sehingga $I \subseteq J_x$.
\end{prop}

\begin{prop} Misalkan $x$ dan $y$ titik-titik dalam ruang Hausdorff kompak. Jika $J_x \subseteq J_y$,
maka $x = y$.
\end{prop}

\begin{prop}\label{0006836} Misalkan $X$ ruang Hausdorff kompak. Subhimpunan $I$ dari $\fml C(X)$
merupakan ideal maksimal dalam $\fml C(X)$ jika dan hanya jika $I = J_x$ untuk suatu $x~\in~X$.
\end{prop}

\begin{cor} Jika $X$ ruang Hausdorff kompak, maka pemetaan $x \mapsto J_x$ dari $X$ ke
$\mx\,\fml C(X)$ adalah bijektif.
\end{cor}

Kekompakan merupakan unsur penting dalam proposisi~\ref{0006836}.

\begin{exam} Dalam aljabar Banach $\fml C_b\bigl(\,(0,1)\,\bigr)$ yang terdiri atas fungsi kontinu
terbatas pada interval $(0,1)$ terdapat ideal maksimal $I$ sedemikian sehingga untuk
\emph{tidak ada} titik $x \in (0,1)$ berlaku $I = J_x$. Misalkan $I$ ideal maksimal yang memuat
ideal $S$ dari semua fungsi $f$ dalam $\fml C_b\bigl(\,(0,1)\,\bigr)$ yang untuknya terdapat
lingkungan $U_f$ dari $0$ dalam $\R$ sedemikian sehingga $f(x) = 0$ untuk semua $x \in U_f \cap (0,1)$.
\end{exam}














\section{Ruang Karakter}
\begin{defn} Suatu
 \index{karakter}%
\df{karakter} (atau
 \index{fungsional!linear!multiplikatif}%
 \index{linear!fungsional!multiplikatif}%
\df{fungsional linear multiplikatif tak nol}) pada aljabar $A$ adalah homomorfisme tak nol
dari $A$ ke dalam~$\C$. Himpunan semua karakter pada $A$ dinotasikan
 \index{delta@$\Delta(A)$!ruang karakter dari~$A$}%
dengan~$\Delta A$.
\end{defn}

\begin{prop}\label{00068511} Misalkan $A$ aljabar beridentitas dan $\phi \in \Delta A$. Maka
 \begin{enumerate}[label=\rm{(\alph*)}]
  \item $\phi(\vc 1) = 1$;
  \item jika $a \in \inv A$, maka $\phi(a) \ne 0$;
  \item jika $a$
 \index{nilpoten}%
  \df{nilpoten} (yaitu, jika $a^n = 0$ untuk suatu $n \in \N$), maka $\phi(a) = 0$;
 \index{idempoten}%
  \item jika $a$ adalah
  \df{idempoten} (yaitu, jika $a^2 = a$), maka $\phi(a)$ adalah $0$ atau~$1$; dan
  \item $\phi(a) \in \sigma(a)$ untuk setiap $a \in A$.
 \end{enumerate}
\end{prop}
Kita catat sambil lalu bahwa bagian (e) dari proposisi sebelumnya tidak memberi kita cara mudah
untuk menunjukkan bahwa spektrum $\sigma(a)$ dari suatu elemen aljabar tidak kosong. Hal ini
bergantung pada pengetahuan bahwa $\Delta(A)$ tidak kosong.

\begin{exam} Pemetaan identitas adalah satu-satunya karakter pada aljabar~$\C$.
\end{exam}

\begin{exam} Misalkan $A$ adalah aljabar matriks berukuran $2 \times 2$, $a = \bigl[a_{ij}\bigr]$,
sedemikian sehingga $a_{12} = 0$. Aljabar ini memiliki tepat dua karakter, yaitu $\phi(a) = a_{11}$
dan $\psi(a) = a_{22}$.
\end{exam}

\begin{proof}[\emph{Petunjuk pembuktian}] Tuliskan $\begin{bmatrix} a & 0 \\ b & c \end{bmatrix}$ sebagai kombinasi linear
dari matriks-matriks $ u = \begin{bmatrix} 1 & 0 \\ 0 & 1 \end{bmatrix}$, $v = \begin{bmatrix} 0 & 0 \\ 1 & 0 \end{bmatrix}$,
dan \smash[t]{$w = \begin{bmatrix} 0 & 0 \\ 0 & 1 \end{bmatrix}$}. Gunakan proposisi~\ref{00068511}.   \ns
\end{proof}

\begin{exam} Aljabar semua matriks $2 \times 2$ dengan entri bilangan kompleks tidak memiliki karakter.
\end{exam}

\begin{prop} Misalkan $A$ aljabar beridentitas dan $\phi$ fungsional linear pada~$A$. Maka $\phi
\in \Delta A$ jika dan hanya jika $\ker\phi$ tertutup terhadap perkalian dan $\phi(\vc 1) =
1$.
\end{prop}

\begin{proof}[\emph{Petunjuk pembuktian}] Untuk arah sebaliknya, terapkan $\phi$ pada hasil kali
$a - \phi(a)\vc 1$ dan $b - \phi(b)\vc 1$ untuk $a$, $b \in A$. \ns
\end{proof}

\begin{prop}\label{00068522} Setiap fungsional linear multiplikatif pada aljabar Banach beridentitas
$A$ kontinu. Bahkan, jika $\phi \in \Delta(A)$, maka $\phi$ merupakan kontraksi dan
$\norm\phi = 1$.
\end{prop}

\begin{exam}\label{00068524} Misalkan $X$ ruang topologis dan $x \in X$. Kita mendefinisikan
 \index{evaluasi!fungsional}%
 \index{evaluasi@$\sbsb EXx$ atau $E_x$ (fungsional evaluasi di~$x$)}%
\df{fungsional evaluasi di $x$}, yang dinotasikan dengan $\sbsb EXx$, melalui
   \[ \sbsb EXx\colon \fml C(X) \sto \C\colon f \mapsto f(x)\,. \]
Fungsional ini merupakan karakter pada~$\fml C(X)$ dan kernelnya adalah $J_x$. Jika hanya ada
satu ruang topologis yang sedang dibahas, kita sederhanakan notasi $\sbsb EXx$ menjadi~$E_x$.
Jadi, khususnya, untuk $f \in \fml C(X)$ kita sering menulis $E_x(f)$ alih-alih
$\sbsb EXx(f)$ yang lebih panjang.
\end{exam}

\begin{defn} Dalam proposisi~\ref{00068522} kita menemukan bahwa setiap karakter pada aljabar Banach
beridentitas $A$ berada pada bola satuan dual~$A^*$. Karena itu kita dapat memberikan himpunan
himpunan $\Delta A$ karakter pada~$A$ dengan topologi-$w^*$ relatif yang diwarisinya
dari~$A^*$. Topologi ini disebut
 \index{Gelfand!topologi}%
 \index{topologi!Gelfand}%
\df{topologi Gelfand} pada $\Delta A$, dan ruang topologis yang dihasilkan disebut
 \index{ruang!karakter}%
 \index{karakter!ruang}%
 \index{ruang!struktur}%
 \index{struktur!ruang}%
\df{ruang karakter} (atau \emph{ruang pembawa}, atau \emph{ruang struktur}, atau \emph{ruang ideal maksimal}, atau \emph{spektrum}) dari~$A$.
\end{defn}

\begin{prop} Ruang karakter dari aljabar Banach beridentitas merupakan ruang Hausdorff kompak.
\end{prop}

\begin{exam}\label{000633} Misalkan $l_1(\Z)$ adalah keluarga semua barisan bilateral
   \[ (\dots,a_{-2},a_{-1},a_0,a_1,a_2,\dots) \]
yang
 \index{mutlak!terjumlahkan!barisan bilateral}%
\df{terjumlahkan mutlak}; yaitu, yang memenuhi $\sum_{k=-\infty}^{\infty} \abs{a_k} <
\infty$. Ini merupakan ruang Banach di bawah operasi penjumlahan dan perkalian skalar titik demi titik,
dengan norma
   \[ \norm a = \sum_{k=-\infty}^{\infty} \abs{a_k}\,. \]
Untuk $a$, $b \in l_1(\Z)$ definisikan $a \ast b$ sebagai barisan yang entri ke-$n$-nya
diberikan oleh
  \[ (a \ast b)_n = \sum_{k=-\infty}^{\infty} a_{n-k}\,b_k\,. \]
Operasi $\ast$ disebut
 \index{konvolusi!dalam $l_1(\Z)$}%
\df{konvolusi}. (Untuk melihat asal definisi ini, cobalah mengalikan deret pangkat
$\sum_{-\infty}^{\infty}a_kz^k$ dan $\sum_{-\infty}^{\infty}b_kz^k$.) Dengan operasi tambahan ini
$l_1(\Z)$ menjadi aljabar Banach komutatif beridentitas.

Ruang ideal maksimal dari aljabar Banach ini
 \index{ruang!ideal maksimal!dari $l_1(\Z)$}%
 \index{ruang karakter!dari $l_1(\Z)$}%
adalah (homeomorfik dengan) lingkaran satuan~$\T$.
\end{exam}

\begin{proof}[\emph{Petunjuk pembuktian}] Untuk setiap $z \in \T$ definisikan
   \[ \psi_z\colon l_1(\Z) \sto \C\colon a \mapsto \sum_{k=-\infty}^\infty a_kz^k\,. \]
Tunjukkan bahwa $\psi_z \in \Delta l_1(\Z)$. Kemudian tunjukkan bahwa pemetaan
   \[ \psi\colon \T \sto \Delta l_1(\Z)\colon z \mapsto \psi_z \]
adalah suatu homeomorfisme. \ns
\end{proof}

\begin{prop} Jika $\phi \in \Delta A$ dengan $A$ aljabar beridentitas, maka $\ker\phi$ merupakan
ideal maksimal dalam~$A$.
\end{prop}

\begin{proof}[\emph{Petunjuk pembuktian}] Untuk menunjukkan kemaksimalan, misalkan $I$ ideal dalam $A$
yang memuat $\ker\phi$ secara sejati. Pilih $z \in I \setminus \ker\phi$. Pertimbangkan elemen
$x - \bigl(\phi(x)/\phi(z)\bigr)\,z$, dengan $x$ elemen sebarang dari~$A$.
\ns
\end{proof}

\begin{prop} Suatu karakter pada aljabar beridentitas ditentukan sepenuhnya oleh kernelnya.
\end{prop}

\begin{proof}[\emph{Petunjuk pembuktian}] Misalkan $a$ elemen aljabar dan $\phi$ suatu karakter.
Untuk berapa banyak bilangan kompleks $\lambda$ dapat $a^2 - \lambda a$ menjadi anggota kernel~$\phi$? \ns
\end{proof}

\begin{cor} Jika $A$ aljabar beridentitas, maka pemetaan $\phi \mapsto \ker\phi$ dari
$\Delta A$ ke~$\mx A$ injektif.
\end{cor}

\begin{prop}\label{0007207} Misalkan $I$ ideal maksimal dalam aljabar Banach komutatif beridentitas~$A$.
Maka terdapat karakter pada $A$ yang kernelnya adalah~$I$.
\end{prop}

\begin{proof}[\emph{Petunjuk pembuktian}] Gunakan korolari~\ref{C073147}. Mengapa kita dapat memandang
pemetaan hasil bagi sebagai suatu karakter?  \ns
\end{proof}

\begin{cor}\label{0007207a} Jika $A$ aljabar Banach komutatif beridentitas, maka pemetaan
$\phi \mapsto \ker\phi$ merupakan bijeksi dari $\Delta A$ ke~$\mx A$.
\end{cor}

\begin{defn} Misalkan $A$ aljabar Banach komutatif beridentitas. Berdasarkan korolari sebelumnya,
kita dapat memberikan $\mx A$ suatu topologi yang membuatnya homeomorfik dengan ruang karakter~$\Delta A$.
Inilah
 \index{ruang!ideal maksimal}%
 \index{ruang!maksimal ideal}%
 \index{maksimal@$\mx A$!homeomorfik dengan $\Delta(A)$}%
\df{ruang ideal maksimal} dari~$A$. Karena $\Delta A$ dan $\mx A$ homeomorfik, lazim untuk
mengidentifikasi keduanya, sehingga $\Delta A$ sering disebut \emph{ruang ideal maksimal} dari~$A$.
\end{defn}

\begin{defn}\label{00073} Misalkan $X$ ruang Hausdorff kompak dan $x \in X$.
Ingat bahwa dalam contoh~\ref{00068524} kita mendefinisikan $\sbsb EXx$, yaitu \emph{fungsional
evaluasi} di~$x$, melalui
   \[ \sbsb EXx(f) := f(x) \]
untuk setiap $f \in \fml C(X)$. Pemetaan
   \[ \sbsb EX \colon X \sto \Delta\fml C(X)\colon x \mapsto \sbsb EXx \]
adalah
 \index{pemetaan evaluasi}%
 \index{evaluasi@$\sbsb EX$!pemetaan evaluasi}%
\df{pemetaan evaluasi} pada~$X$. Sebagaimana disebutkan sebelumnya, ketika hanya satu ruang topologis
yang sedang dipertimbangkan, biasanya kita menyingkat $\sbsb EX$ menjadi $E$ dan $\sbsb EXx$ menjadi~$E_x$.
\end{defn}

\begin{notn} Untuk menyatakan bahwa dua ruang topologis $X$ dan $Y$ homeomorfik, kita
 \index{<binrel@$X \approx Y$ ($X$ dan $Y$ homeomorfik)}%
menulis $X \approx Y$.
\end{notn}

\begin{prop}\label{000731} Misalkan $X$ ruang Hausdorff kompak. Maka
 \index{ruang!ideal maksimal!dari $\fml C(X)$}%
 \index{ruang karakter!dari $\fml C(X)$}%
pemetaan evaluasi pada~$X$
   \[ \sbsb EX \colon X \sto \Delta\fml C(X) \colon x \mapsto \sbsb EXx \]
merupakan homeomorfisme. Jadi kita memiliki
   \[ X \approx \Delta\fml C(X) \approx \mx\fml C(X)\,. \]
Lebih dari itu: bukan hanya setiap $\sbsb EX$ merupakan homeomorfisme antara ruang Hausdorff kompak,
melainkan $E$ sendiri merupakan suatu \emph{ekuivalensi natural} antara funktor---funktor identitas
dan funktor $\Delta\fml C$.
\end{prop}

Identifikasi antara ruang Hausdorff kompak $X$ dan ruang karakternya serta ruang ideal maksimalnya
begitu kuat sehingga banyak orang membicarakannya seolah-olah ketiganya benar-benar sama. Sangat lazim
untuk mendengar, misalnya, bahwa ``ideal maksimal dalam $\fml C(X)$ hanyalah titik-titik di~$X$''.
Walaupun secara harfiah tidak benar, ungkapan itu memang terdengar sedikit kurang mengintimidasi daripada
``ideal maksimal dari $\fml C(X)$ adalah kernel fungsional evaluasi pada titik-titik~$X$''.

\begin{prop}\label{000735} Misalkan $X$ dan $Y$ ruang Hausdorff kompak dan
 \index{funktor!$\fml C$ sebagai suatu}%
 \index{C@$\fml C$ (funktor)}%
$F \colon X \sto Y$ kontinu. Ingat bahwa dalam contoh~\ref{C029717} kita mendefinisikan $\fml C(F)$
pada $\fml C(Y)$ melalui
   \[ \fml C(F)(g) = g \circ F\]
untuk semua $g \in \fml C(Y)$. Maka
  \begin{enumerate}[label=\rm{(\alph*)}]
   \item $\fml C(F)$ memetakan $\fml C(Y)$ ke dalam $\fml C(X)$.
   \item Pemetaan $\fml C(F)$ merupakan homomorfisme aljabar Banach beridentitas yang kontraktif.
   \item $\fml C(F)$ injektif jika dan hanya jika $F$ surjektif.
   \item $\fml C(F)$ surjektif jika dan hanya jika $F$ injektif.
   \item Jika $X$ homeomorfik dengan $Y$, maka $\fml C(X)$ isomorfik secara isometrik
   dengan~$\fml C(Y)$.
  \end{enumerate}
\end{prop}

\begin{prop}\label{000736} Misalkan $A$ dan $B$ aljabar Banach komutatif beridentitas dan
 \index{funktor!$\Delta$ sebagai suatu}%
 \index{delta@$\Delta$ (funktor)}%
$T\colon A \sto B$ homomorfisme aljabar beridentitas. Definisikan $\Delta T$ pada $\Delta B$ melalui
   \[ \Delta T(\psi) = \psi \circ T \]
untuk semua $\psi \in \Delta B$. Maka
  \begin{enumerate}[label=\rm{(\alph*)}]
   \item $\Delta T$ memetakan $\Delta B$ ke dalam $\Delta A$.
   \item Pemetaan $\Delta T\colon \Delta B \sto \Delta A$ kontinu.
   \item Jika $T$ surjektif, maka $\Delta T$ injektif.
   \item Jika $T$ merupakan suatu isomorfisme (aljabar), maka $\Delta T$ merupakan homeomorfisme.
   \item Jika $A$ dan $B$ isomorfik (secara aljabar), maka $\Delta A$ dan $\Delta B$
   homeomorfik.
  \end{enumerate}
\end{prop}

\begin{cor} Misalkan $X$ dan $Y$ ruang Hausdorff kompak. Jika $\fml C(X)$ dan $\fml C(Y)$ isomorfik
secara aljabar, maka $X$ dan $Y$ homeomorfik.
\end{cor}

\begin{cor}\label{nasc_homeomor_CHSps} Dua ruang Hausdorff kompak homeomorfik jika dan hanya jika aljabar
fungsi kontinu bernilai kompleksnya isomorfik (secara aljabar).
\end{cor}

Korolari~\ref{nasc_homeomor_CHSps} menimbulkan sekumpulan masalah yang sangat menarik. Korolari itu menyarankan
bahwa untuk setiap sifat topologis ruang Hausdorff kompak $X$ seharusnya ada sifat aljabar yang bersesuaian
dari aljabar $\fml C(X)$, dan \emph{sebaliknya}. Namun, korolari tersebut tidak memberikan resep yang dapat
digunakan untuk menemukan korespondensi ini. Salah satu contoh yang sangat sederhana dari hasil semacam ini
diberikan di bawah sebagai proposisi~\ref{prop_connected_idempotent}. Untuk pembahasan yang indah mengenai
pertanyaan ini, lihat monograf Semadeni~\cite{Semadeni:1971}, khususnya tabel pada halaman~132--133.

\begin{cor} Misalkan $X$ dan $Y$ ruang Hausdorff kompak. Jika $\fml C(X)$ dan $\fml C(Y)$ isomorfik
secara aljabar, maka keduanya isomorfik secara isometrik.
\end{cor}

\begin{prop}\label{prop_connected_idempotent} Ruang topologis $X$ terhubung jika dan hanya jika aljabar $\fml C(X)$
tidak memuat idempoten nontrivial.
\end{prop}

(Idempoten \emph{trivial} suatu aljabar adalah $\vc 0$ dan~$\vc 1$.)













\section{Transformasi Gelfand}
\begin{defn} Misalkan $A$ aljabar Banach komutatif dan $a\in A$. Definisikan
   \[\sbsb{\Gamma}A a\colon \Delta A \sto \C \colon \phi \mapsto \phi(a) \]
untuk setiap $\phi \in \Delta(A)$. (Notasi alternatif: jika tidak menimbulkan kebingungan, kita
sering menulis $\Gamma a$ atau $\wh a$ untuk $\sbsb{\Gamma}A a$.) Pemetaan $\sbsb{\Gamma}A$ adalah
 \index{Gelfand!transformasi!pada aljabar Banach}%
 \index{Gamma@$\sbsb\Gamma A$, $\Gamma$ (transformasi Gelfand)}%
 \index{<unoptop@$\wh a$ (transformasi Gelfand dari~$a$)}%
\df{transformasi Gelfand pada~$A$}.
\end{defn}

Karena $\Delta A \subseteq A^*$, jelas bahwa $\sbsb{\Gamma}A a$ hanyalah restriksi $a^{**}$
pada ruang karakter~$A$. Lebih lanjut, topologi Gelfand pada $\Delta A$ adalah topologi-$w^*$ relatif,
yaitu topologi terlemah yang membuat $a^{**}$ kontinu pada $\Delta A$ untuk setiap $a \in A$,
sehingga $\sbsb{\Gamma}A a$ merupakan fungsi kontinu pada $\Delta A$. Jadi
$\sbsb{\Gamma}A \colon A \sto \fml C(\Delta A)$.

Sebagai ringkasan dan demi kemudahan, elemen $\sbsb{\Gamma}A a = \Gamma a = \wh a$ biasanya cukup
disebut
 \index{Gelfand!transformasi!dari suatu elemen}%
\emph{transformasi Gelfand dari~$a$}---karena frasa \emph{transformasi Gelfand pada $A$ yang dievaluasi
pada~$a$} terasa janggal.

\begin{defn} Kita mengatakan bahwa suatu keluarga fungsi $\fml F$ pada himpunan $S$
 \index{pemisahan!titik}%
\df{memisahkan titik-titik} pada $S$ (atau merupakan
 \index{pemisah!keluarga fungsi}%
\df{keluarga fungsi pemisah} pada $S$) jika untuk setiap pasangan elemen berbeda $x$
dan $y$ dari $S$ terdapat $f \in \fml F$ sedemikian sehingga $f(x) \ne f(y)$.
\end{defn}

\begin{prop} Misalkan $X$ ruang topologis kompak. Maka $\fml C(X)$ memisahkan titik-titik jika dan
hanya jika $X$ Hausdorff.
\end{prop}

\begin{prop}\label{0007402} Jika $A$ aljabar Banach komutatif, maka $\sbsb{\Gamma}A
\colon A \sto \fml C(\Delta A)$ merupakan homomorfisme aljabar yang kontraktif dan jangkauan
$\sbsb{\Gamma}A$ merupakan subaljabar pemisah dari $\fml C(\Delta A)$. Jika, sebagai tambahan,
$A$ beridentitas, maka $\Gamma$ bernorma satu, jangkauannya merupakan subaljabar beridentitas dari
$\fml C(\Delta A)$, dan pemetaan itu merupakan homomorfisme beridentitas.
\end{prop}

\begin{prop} Misalkan $a$ elemen aljabar Banach komutatif beridentitas~$A$. Maka $a$ invertibel dalam
$A$ jika dan hanya jika $\hat a$ invertibel dalam~$\fml C(\Delta A)$.
\end{prop}

\begin{prop}\label{00074043} Misalkan $A$ aljabar Banach komutatif beridentitas dan $a$ suatu
elemen dari~$A$. Maka $\ran\hat a = \sigma(a)$ dan $\norm{\hat a}_u~=~\rho(a)$.
\end{prop}

\begin{defn} Suatu elemen $a$ dari aljabar Banach disebut
 \index{kuasinilpoten}%
\df{kuasinilpoten} jika $\lim\limits_{n \sto \infty}\norm{a^n}^{1/n} = 0$.
\end{defn}

\begin{prop} Misalkan $a$ elemen aljabar Banach komutatif beridentitas~$A$. Maka pernyataan-pernyataan
berikut ekuivalen:
  \begin{enumerate}[label=\rm{(\alph*)}]
   \item $a$ kuasinilpoten;
   \item $\rho(a) = 0$;
   \item $\sigma(a) = \{0\}$;
   \item $\Gamma a = 0$;
   \item $\phi(a) = 0$ untuk setiap $\phi \in \Delta A$;
   \item $a \in \bigcap\mx A$.
  \end{enumerate}
\end{prop}

\begin{defn} Suatu aljabar Banach disebut
 \index{semisederhana}%
\df{semisederhana} jika tidak memiliki elemen kuasinilpoten tak nol.
\end{defn}

\begin{prop} Misalkan $A$ aljabar Banach komutatif beridentitas. Pernyataan-pernyataan berikut ekuivalen:
  \begin{enumerate}[label=\rm{(\alph*)}]
   \item $A$ semisederhana;
   \item jika $\rho(a) = 0$, maka $a = 0$;
   \item jika $\sigma(a) = \{0\}$, maka $a = 0$;
   \index{monomorfisme}%
   \item transformasi Gelfand $\sbsb{\Gamma}A$ merupakan monomorfisme (yaitu,
   homomorfisme injektif);
   \item jika $\phi(a) = 0$ untuk setiap $\phi \in \Delta A$, maka $a = 0$;
   \item $\bigcap\mx A = \{0\}$.
  \end{enumerate}
\end{prop}

\begin{prop}\label{00074083} Misalkan $A$ aljabar Banach komutatif beridentitas. Maka
pernyataan-pernyataan berikut ekuivalen:
  \begin{enumerate}[label=\rm{(\alph*)}]
    \item $\|a^2\| = \|a\|^2$ untuk semua $a \in A$;
    \item $\rho(a) = \|a\|$ untuk semua $a \in A$; dan
    \item transformasi Gelfand merupakan isometri; yaitu, $\norm{\wh a}_u = \|a\|$ untuk
    semua $a \in A$.
  \end{enumerate}
\end{prop}

\begin{exam} Transformasi Gelfand pada $l_1(\Z)$ bukan isometri.
\end{exam}

Ingat dari proposisi~\ref{000731} bahwa ketika $X$ ruang Hausdorff kompak, pemetaan evaluasi $E_X$
mengidentifikasi ruang $X$ dengan ruang ideal maksimal aljabar Banach~$\fml C(X)$. Jadi menurut proposisi~\ref{000735},
pemetaan $\fml CE_X$ mengidentifikasi aljabar $\fml C(\Delta(\fml C(X)))$ fungsi kontinu pada ruang ideal maksimal
ini dengan aljabar $\fml C(X)$ itu sendiri. Ternyata transformasi Gelfand pada aljabar $\fml C(X)$ hanyalah invers
dari pemetaan identifikasi ini.

\begin{exam} Misalkan $X$ ruang Hausdorff kompak. Maka transformasi Gelfand pada aljabar Banach
$\fml C(X)$ merupakan isomorfisme isometrik. Bahkan, pada $\fml C(X)$ transformasi Gelfand
\smash[b]{$\sbsb\Gamma X$} adalah $(\fml CE_X)^{-1}$.
\end{exam}

\begin{exam}\label{0007472} Ruang ideal maksimal aljabar Banach $L_1(\R)$
 \index{ruang!ideal maksimal!untuk $L_1(\R)$}%
adalah $\R$ sendiri dan transformasi Gelfand pada $L_1(\R)$ adalah transformasi Fourier.
\end{exam}

\begin{proof} Lihat~\cite{BaggettF:1979}, halaman 169--171.  \ns
\end{proof}

Fungsi $t \mapsto e^{it}$ merupakan bijeksi dari interval $[-\pi,\pi)$ ke
 \index{T@$\T$ (lingkaran satuan pada bidang kompleks)}%
lingkaran satuan~$\T$ pada bidang kompleks. Salah satu akibatnya adalah kita tidak perlu membedakan antara
  \begin{enumerate}[label=\rm{(\alph*)}]
     \item fungsi $2\pi$-periodik pada~$\R$,
     \item semua fungsi pada $[-\pi,\pi)$,
     \item fungsi $f$ pada $[-\pi,\pi]$ sedemikian sehingga $f(-\pi) = f(\pi)$, dan
     \item fungsi pada~$\T$.
  \end{enumerate}
Dalam kelanjutan, kita akan sering mengidentifikasi kelas-kelas fungsi ini tanpa penjelasan lebih lanjut.

Identifikasi lain yang berguna adalah antara lingkaran satuan $\T$ dalam $\C$ dan ruang ideal maksimal
aljabar $l_1(\Z)$. Homeomorfisme $\psi$ antara dua ruang Hausdorff kompak ini didefinisikan dalam contoh~\ref{000633}.
Dalam bekerja dengan transformasi Gelfand $\Gamma_a$ dari elemen $a \in l_1(\Z)$, sering secara teknis lebih mudah
memperlakukannya sebagai fungsi, sebut saja $G_a$, yang daerah asalnya adalah~$\T$, seperti ditunjukkan oleh diagram berikut.
  \[
    \xy
      \btriangle[\T`\Delta l_1(\Z)`\C;\psi`G_a`\Gamma_a]
    \endxy
  \]
Jadi untuk $a \in l_1(\Z)$ dan $z \in \T$ berlaku
  \[ G_a(z) = \Gamma_a(\psi_z) = \psi_z(a) = \sum_{k=-\infty}^\infty a_kz^k\,. \]

\begin{defn} Jika $f \in \fml L_1([-\pi,\pi))$, maka
  \index{Fourier!deret}%
\df{deret Fourier} untuk $f$ adalah deret
   \[ \sum_{n=-\infty}^\infty\wt f(n)\exp(int) \qquad -\pi \le t \le \pi \]
dengan barisan $\wt f$ didefinisikan oleh
   \[ \wt f(n) = {\textstyle\frac1{2\pi}}\int_{-\pi}^\pi f(t)\exp(-int)\,dt \]
untuk semua $n \in \Z$. Barisan dua arah tak hingga $\wt f$ adalah
  \index{Fourier!transformasi!pada $L_1([-\pi,\pi])$}%
\df{transformasi Fourier} dari $f$, dan bilangan $\wt f(n)$ adalah
  \index{Fourier!koefisien}%
\df{koefisien Fourier} ke-$n$ dari $f$. Jika $\wt f \in l_1(\Z)$, kita mengatakan bahwa $f$ memiliki
  \index{mutlak!konvergen!deret Fourier}%
  \index{Fourier!deret!konvergen mutlak}%
\df{deret Fourier yang konvergen mutlak}. Himpunan semua fungsi kontinu pada $\T$
dengan deret Fourier yang konvergen mutlak dinotasikan dengan
  \index{mutlak@$\fml{A\,C}(\T)$ (deret Fourier yang konvergen mutlak}%
$\fml{A\,C}(\T)$.
\end{defn}

\begin{prop} Jika $f$ fungsi kontinu $2\pi$-periodik pada $\R$ yang transformasi Fouriernya nol,
maka $f = 0$.
\end{prop}

\begin{cor} Transformasi Fourier pada $\fml C(\T)$ injektif.
\end{cor}

\begin{prop} Transformasi Fourier pada $\fml C(\T)$ merupakan invers kiri dari transformasi Gelfand
pada~$l_1(\Z)$.
\end{prop}

\begin{prop} Jangkauan transformasi Gelfand pada $l_1(\Z)$ adalah aljabar Banach komutatif beridentitas
$\fml{A\,C}(\T)$.
\end{prop}

\begin{prop} Terdapat fungsi kontinu yang deret Fouriernya divergen di~$0$.
\end{prop}

\begin{proof} Lihat, misalnya, \cite{HewittS:1965}, latihan 18.45.  \ns
\end{proof}

Apa yang dikatakan hasil sebelumnya mengenai transformasi Gelfand $\Gamma\colon l_1(\Z) \sto
\fml C(\T)$?

\vskip 3 pt

Misalkan suatu fungsi $f$ termasuk dalam $\fml{AC}(\T)$ dan tidak pernah nol. Maka $1/f$ tentu kontinu pada~$\T$,
tetapi apakah deret Fouriernya konvergen mutlak? Salah satu keberhasilan awal studi abstrak aljabar Banach
adalah pembuktian yang sangat sederhana atas pertanyaan ini, yang mula-mula diberikan oleh Norbert Wiener.
Bukti asli Wiener melalui perbandingan cukup sulit.

\begin{thm}[Teorema Wiener] Misalkan $f$ fungsi kontinu pada $\T$ yang tidak pernah
 \index{teorema Wiener}%
nol. Jika $f$ memiliki deret Fourier yang konvergen mutlak, maka kebalikannya~$1/f$ juga demikian.
\end{thm}

\begin{exam} Transformasi Laplace juga dapat dipandang sebagai kasus khusus transformasi Gelfand.
Untuk rinciannya, lihat~\cite{BaggettF:1979}, halaman 173--175.
\end{exam}













\section{Aljabar-$C^*$ Beridentitas}

\begin{prop} Jika $a$ elemen suatu aljabar-$*\,$, maka $\sigma(a^*) = \conj{\sigma(a)}$.
\end{prop}


\begin{prop}\label{001311} Dalam setiap aljabar-$C^*$, involusi merupakan isometri. Artinya,
 \index{involusi!merupakan isometri}%
$\norm{a^*} = \norm{a}$ untuk setiap elemen $a$ dalam aljabar.
\end{prop}

Dalam definisi~\ref{def_norm_alg} tentang \emph{aljabar bernorma}, kita memberikan syarat khusus bahwa
elemen identitas dari aljabar bernorma beridentitas harus bernorma satu. Dalam aljabar-$C^*$, syarat
ini bersifat berlebihan.

\begin{cor} Dalam aljabar-$C^*$ beridentitas, $\norm{\vc 1} = 1$.
\end{cor}

\begin{cor}\label{001311005fa} Setiap elemen uniter dalam aljabar-$C^*$ beridentitas bernorma satu.
\end{cor}

\begin{cor}\label{001311008} Jika $a$ elemen aljabar-$C^*$ $A$ sedemikian sehingga
$ab = \vc 0$ untuk setiap $b \in A$, maka $a = \vc 0$.
\end{cor}

\begin{prop}\label{00131101} Misalkan $a$ elemen normal dari aljabar-$C^*$.
Maka $\norm{a^2} = {\norm a}^2$. Jika aljabar tersebut beridentitas, maka $\rho(a)~=~\norm a$.
\end{prop}

\begin{cor} Jika $p$ proyeksi tak nol dalam aljabar-$C^*$, maka $\norm p = 1$.
\end{cor}

\begin{cor}\label{00131102} Misalkan $a \in A$, dengan $A$ suatu aljabar-$C^*$ komutatif. Maka
$\norm{a^2} = {\norm a}^2$ dan, jika sebagai tambahan $A$ beridentitas, maka $\rho(a)~=~\norm a$.
\end{cor}

\begin{cor}\label{00131103} Pada aljabar-$C^*$ komutatif beridentitas $A$, transformasi Gelfand
$\Gamma$ merupakan isometri; yaitu, $\norm{\Gamma_a}_u = \norm{\wh a}_u = \norm a$ untuk
setiap $a \in A$.
\end{cor}

\begin{cor}\label{00131104} Norma aljabar-$C^*$ beridentitas bersifat unik, dalam arti bahwa
untuk aljabar beridentitas $A$ dengan involusi, terdapat paling banyak satu norma yang menjadikan
$A$ sebuah aljabar-$C^*$.
\end{cor}

\begin{prop}\label{001312} Jika $a$ elemen swaadjoin dari aljabar-$C^*$ beridentitas,
 \index{spektrum!elemen swaadjoin}%
 \index{swaadjoin!spektrum elemen}%
maka $\sigma(a) \subseteq \R$.
\end{prop}

\begin{proof}[\emph{Petunjuk pembuktian}] Tuliskan $\lambda \in \sigma(a)$ dalam bentuk $\alpha + i\beta$, dengan $\alpha$, $\beta \in \R$.
Untuk setiap $n \in \N$ misalkan $b_n := a - \alpha \vc 1 + i n\beta\vc 1$. Tunjukkan bahwa untuk setiap $n$ skalar
$i(n+1)\beta$ termasuk dalam spektrum $b_n$ dan karena itu
   \[ {\abs{i(n+1)\beta}}^2 \le {\norm{a - \alpha\vc 1}}^2 + n^2\beta^2\,. \]   \ns
\end{proof}

\begin{cor}\label{001406} Misalkan $A$ aljabar-$C^*$ komutatif beridentitas. Jika $a \in A$ swaadjoin,
maka transformasi Gelfand $\wh a$ bernilai real.
\end{cor}

\begin{cor}\label{0014062} Setiap karakter pada aljabar-$C^*$ merupakan homomorfisme-$*\,$.
\end{cor}

\begin{cor}\label{0014062fa} Transformasi Gelfand pada aljabar-$C^*$ komutatif beridentitas $A$ merupakan
homomorfisme-$*\,$ beridentitas ke dalam~$\fml C(\Delta A)$.
\end{cor}

\begin{cor} Jangkauan transformasi Gelfand pada aljabar-$C^*$ komutatif beridentitas $A$ merupakan
subaljabar pemisah, swaadjoin, beridentitas dari~$\fml C(\Delta A)$.
\end{cor}

\begin{prop}\label{001313} Jika $u$ elemen uniter dari aljabar-$C^*$ beridentitas,
 \index{spektrum!elemen uniter}%
 \index{uniter!spektrum elemen}%
maka $\sigma(u) \subseteq \T$.
\end{prop}

\begin{prop} Transformasi Gelfand $\Gamma$ merupakan ekuivalensi natural antara funktor identitas
dan funktor $\fml C\Delta$ pada kategori aljabar-$C^*$ komutatif beridentitas dan homomorfisme-$*\,$
beridentitas.
\end{prop}













\section{Teorema Gelfand-Naimark}

Kita mengingat kembali suatu teorema baku dari analisis real.

\begin{thm}[Teorema Stone-Weierstrass]\label{00141} Misalkan $X$ ruang Hausdorff kompak. Setiap
 \index{teorema Stone-Weierstrass}%
subaljabar-$*\,$ pemisah beridentitas dari $\fml C(X)$ rapat.
\end{thm}

\begin{proof} Lihat \cite{Douglas:1972}, teorema 2.40. \ns \end{proof}

Versi pertama dari \emph{teorema Gelfand-Naimark} menyatakan bahwa setiap aljabar-$C^*$ komutatif
beridentitas adalah aljabar semua fungsi kontinu pada suatu ruang Hausdorff kompak. Tentu saja kata
\emph{adalah} dalam kalimat sebelumnya berarti \emph{isomorfik secara isometrik dan sebagai-$*$ dengan}.
Ruang Hausdorff kompak yang dimaksud adalah ruang karakter aljabar tersebut.

\begin{thm}[Teorema Gelfand-Naimark I]\label{00142} Misalkan $A$ aljabar-$C^*$ komutatif
 \index{Gelfand-Naimark!teorema I}%
beridentitas. Maka transformasi Gelfand $\Gamma_A\colon a \mapsto \wh a$ merupakan
isomorfisme-$*\,$ isometrik dari $A$ ke~$\fml C(\Delta A)$.
\end{thm}






\begin{defn} Misalkan $S$ subhimpunan tak kosong dari aljabar-$C^*$~$A$. Irisan
keluarga semua subaljabar-$C^*$ dari $A$ yang memuat $S$ disebut
 \index{C*@$C^*$-aljabar!yang dibangkitkan oleh suatu himpunan}%
\df{subaljabar-$C^*$ yang dibangkitkan oleh~$S$}. Kita menotasikannya dengan
 \index{C*@$C^*(S)$ ($C^*$-subalgebra yang dibangkitkan oleh~$S$)}%
$C^*(S)$. (Mudah dilihat bahwa irisan suatu keluarga subaljabar-$C^*$ memang merupakan
aljabar-$C^*$.) Dalam beberapa kasus kita sedikit menyingkat notasi: misalnya, jika $a \in A$ kita
menulis $C^*(a)$ untuk~$C^*(\{a\})$ dan $C^*(\vc 1,a)$ untuk $C^*(\{\vc 1,a\})$.
\end{defn}

\begin{prop} Misalkan $S$ subhimpunan tak kosong dari aljabar-$C^*$~$A$. Untuk setiap bilangan natural
$n$ definisikan himpunan $W_n$ sebagai himpunan semua elemen $a$ dari $A$ yang untuknya terdapat
$x_1$, $x_2$, \dots, $x_n$ dalam $S \cup S^*$ sedemikian sehingga $a = x_1x_2\cdots x_n$. Misalkan
$W = \bigcup_{n=1}^\infty W_n$. Maka
   \[ C^*(S) = \clo{\spn W}\,. \]
\end{prop}













\begin{thm}[Teorema Spektral Abstrak]\label{00144} Jika $a$ elemen normal dari aljabar-$C^*$
 \index{abstrak!teorema spektral}%
 \index{spektral!teorema!abstrak}%
beridentitas $A$, maka aljabar-$C^*$ $\fml C(\sigma(a))$ isomorfik secara isometrik dan sebagai-$*$
dengan $C^*(\vc 1,a)$.
\end{thm}

\begin{proof}[\emph{Petunjuk pembuktian}] Gunakan transformasi Gelfand dari $a$ untuk mengidentifikasi
ruang ideal maksimal $C^*(\vc 1,a)$ dengan spektrum~$a$. Terapkan funktor~$\fml C$.
Komposisikan pemetaan yang dihasilkan dengan $\Gamma^{-1}$, dengan $\Gamma$ transformasi Gelfand pada
aljabar-$C^*$~$C^*(\vc 1,a)$.   \ns
\end{proof}

Jika $\psi\colon \fml C(\sigma(a)) \sto C^*(\vc 1,a)$ adalah isomorfisme-$*\,$ pada teorema sebelumnya dan $f$
merupakan fungsi kontinu pada spektrum $a$, maka elemen $\psi(f)$ dalam aljabar $A$ biasanya dinotasikan
dengan~$f(a)$. Operasi yang mengaitkan fungsi kontinu $f$ dengan elemen $f(a)$ dalam $A$ ini disebut
 \index{kalkulus!fungsional}%
 \index{fungsional!kalkulus}%
\emph{kalkulus fungsional} yang terkait dengan~$a$.

\begin{exam} Misalkan pada teorema sebelumnya $\psi\colon \fml C(\sigma(a)) \sto
C^*(\vc 1,a)$ merealisasikan isomorfisme-$*\,$ isometrik. Maka citra di bawah $\psi$ dari fungsi konstan~$\vc 1$
pada spektrum $a$ adalah $\vc 1_A$, dan citra di bawah $\psi$ dari fungsi identitas $\lambda \mapsto \lambda$
pada spektrum $a$ adalah~$a$.
\end{exam}

\begin{exam}\label{X_sqroot_op} Misalkan $T$ operator normal pada ruang Hilbert $H$ yang spektrumnya termuat
dalam~$[0,\infty)$. Misalkan $\psi\colon \fml C(\sigma(T)) \sto C^*(I,T)$ merealisasikan isomorfisme-$*\,$
isometrik antara kedua aljabar-$C^*$ ini. Maka terdapat setidaknya satu operator $\sqrt T$ yang kuadratnya
adalah~$T$. Memang, setiap kali $f$ fungsi kontinu pada spektrum operator normal~$T$, kita dapat secara bermakna
membicarakan operator~$f(T)$.
\end{exam}

\begin{exam} Jika $N$ operator normal pada ruang Hilbert yang spektrumnya adalah $\{0,1\}$, maka
$N$ merupakan proyeksi ortogonal.
\end{exam}

\begin{exam} Jika $N$ operator normal pada ruang Hilbert yang spektrumnya termuat dalam lingkaran satuan~$\T$,
maka $N$ merupakan operator uniter.
\end{exam}

\begin{prop}[Teorema pemetaan spektral]\label{001446} Misalkan $a$ elemen swaadjoin
 \index{spektral!teorema pemetaan}%
dalam aljabar-$C^*$ beridentitas~$A$ dan $f$ fungsi kontinu bernilai kompleks pada spektrum
$a$. Maka
  \[ \sigma(f(a)) = f^{\sto}(\sigma(a))\,. \]
\end{prop}

Kalkulus fungsional memetakan fungsi kontinu ke fungsi kontinu dalam pengertian berikut.

\begin{prop}\label{001449} Misalkan $A$ aljabar-$C^*$ beridentitas, $K$ subhimpunan kompak tak kosong dari~$\R$,
dan $f\colon K \sto \C$ fungsi kontinu. Misalkan $\ofml H_K$ himpunan semua elemen swaadjoin dari $A$ yang spektrumnya
merupakan subhimpunan dari~$K$. Maka fungsi $f\colon \ofml H_K \sto A\colon h \mapsto f(h)$ kontinu.
\end{prop}

\begin{proof}[\emph{Petunjuk pembuktian}] Aproksimasi $f$ secara seragam dengan suatu polinom dan gunakan
\emph{teorema pemetaan spektral}~\ref{001446} serta argumen $\epsilon$ dibagi $3$.   \ns
\end{proof}













\endinput
